
\section{Introduction}

The universal characterization of algebraic K-theory by
Blumberg, Gepner, and Tabuada provides a systematic approach
to additive invariants of stable infinity-categories
\citep{BGT}. Mazel-Gee and Stern extended this perspective
to stable $(\infty,2)$-categories, obtaining category-valued
additive motives and secondary algebraic K-theory \citep{MGS}.
In this paper, we develop a theory of higher additive motives
that extends these constructions to stable $(\infty,n)$-categories.

Our framework is based on the symmetric monoidal tower
\[
 \nPr0=\Sp,\qquad
 \nPr{(r+1)}=\Mod_{\nPr r}(\PrL),
\]
whose precise definition is given in
Section~\ref{sec:conventions}.
The relevant exact sequences generalize cofiber sequences
and the localization sequences of Mazel-Gee and Stern,
incorporating adjunctions and exactness in the Hom objects.

Our first main result constructs universal additive
invariants at every finite height.

\begin{theorem}[Additive motives]
\label{thm:main}
For every $n\geq1$ and $A\in\nPr{(n+1)}$,
there exist $\M_n(A)\in\nPr n$ and a canonical
additive invariant $U_A:A\to \M_n(A)$ such that
\begin{equation}\label{eq:main-up}
 \Fun^L_{\nPr{(n-1)}}(\M_n(A),E)
 \simeq \Add_{n,\mathrm{filt}}(A,E)
\end{equation}
naturally in $E\in\nPr n$.
\end{theorem}

The right-hand side, $\Add_{n,\mathrm{filt}}(A,E)$, is the category of additive invariants valued in $E$.

We give two constructions of these motives.
The first is obtained by localizing enriched presheaves
(Theorem~\ref{thm:presheaf-universal}),
generalizing the constructions of additive motives in
\cite{BGT,MGS}. The second is a higher analogue of
Waldhausen's $S$-construction
(Theorem~\ref{thm:weighted-s-universal}).
A key ingredient is the existence of compact modules
classifying exact sequences.
We establish the universal properties of both constructions,
thereby proving their equivalence.

Having constructed the motives, we compute their Hom objects.
This reveals a recursive relationship between additive motives
at successive heights.

\begin{theorem}[Mapping objects;
Theorem~\ref{thm:mapping-comparison}]
\label{thm:intro-mapping}
Let $A\in\nPr{(n+1)}$, let $c\in A$ be compact,
and let $a\in A$ be arbitrary. Then
\[
 \operatorname{map}_{\M_1(A)}(U_Ac,U_Aa)
 \simeq K\bigl(\PrHom_A(c,a)^\omega\bigr)
 \qquad(n=1),
\]
and
\[
 \PrHom_{\M_n(A)}(U_Ac,U_Aa)
 \simeq \M_{n-1}\bigl(\PrHom_A(c,a)\bigr)
 \qquad(n\geq2),
\]
where $K$ denotes connective algebraic K-theory.
These equivalences are compatible with composition,
units, and external products.
\end{theorem}

In other words, the higher additive motives themselves
are the values of higher K-theory.

These results recover the additive motives of \cite{BGT,HSS}
over rigid stable bases and the category-valued
secondary motives of Mazel-Gee and Stern.

Finally, we study the interaction of higher K-theory
with Stefanich rings. 
For $A\in\nPr n$, higher K-theory is defined
by
\[
K^{[n]}(A)=
\begin{cases}
\operatorname{map}_{\M_1(\nPr1)}
\bigl(U_{\nPr1}(\Sp),U_{\nPr1}(A)\bigr),
&n=1,\\
\PrHom_{\M_n(\nPr n)}
\bigl(U_{\nPr n}(\nPr{(n-1)}),U_{\nPr n}(A)\bigr),
&n\geq2.
\end{cases}
\]
The multiplicative compatibility of our
mapping formula yields the following result.

\begin{theorem}[K-theory of Stefanich rings;
Theorem~\ref{thm:stefanich-k}]
\label{thm:intro-stefanich}
Higher K-theory defines a functor
\[
 K_{\mathrm{St}}:
 \operatorname{StRing}\longrightarrow
 \operatorname{StRing},
 \qquad
 K_{\mathrm{St}}(\mathbf R)_r
 =K^{[r+1]}(R_{r+1}).
\]
\end{theorem}




The paper is organized as follows.
Section~\ref{sec:conventions} introduces exact sequences
and additive diagrams.
Section~\ref{sec:weights} constructs universal exact sequences.
Section~\ref{sec:presheaf} constructs additive motives by
localizing enriched presheaves, while
Section~\ref{sec:waldhausen} gives an alternative construction
using a weighted Waldhausen $S$-construction on exact flags.
Section~\ref{sec:mapping} establishes the mapping formulas
and comparisons with existing theories.
Section~\ref{sec:higher-k} develops higher K-theory.
Finally, Section~\ref{sec:stefanich} studies the K-theory
of Stefanich rings.
